\section{Method and evidence discipline}

This article is a \emph{critical integrative synthesis and operational
proposal}. It is not a systematic review or a claim of exhaustive retrieval.
The project registry records the query strata, primary-source checks, source
notes, competing explanations, falsification criteria, and a claim--evidence
matrix. Searches were expanded adversarially whenever a field offered a likely
redundancy account. The final bounded closure covered classical
reachability/viability, causal and semantic information, feedback information
thermodynamics, positive and critical Quantum Darwinism, quantum reference
frames, biological organization and carrier interventions, and emergent
quasiparticles. The bounded cross-domain search cutoff was 20 July 2026. A
targeted 21 July amendment added two primary Gravity from Entropy sources and
the formal IOTA--GfE bridge discussed below. A second question-triggered
amendment on 21 July reviewed the adjacent possible, correlated-novelty models,
and the Theory of the Adjacent Possible (TAP). Neither amendment reopened the
other search strata or converted the review into an exhaustive search.

Claims are assigned to five evidence strata:
\begin{enumerate}
  \item \textbf{Established field result}: a theorem, experiment, or formal
  construction within its native assumptions.
  \item \textbf{Project constructive evidence}: a frozen or prospectively
  registered construction within a declared model.
  \item \textbf{Project empirical/computational evidence}: measurements on an
  engineered or observed data-generating system, without reclassifying that
  system as fundamental physics.
  \item \textbf{Synthesis inference}: a cross-field structural mapping that may
  be supported even when no source states it.
  \item \textbf{IOTA/ICP hypothesis}: a proposed universality, emergence, or
  ontology claim.
\end{enumerate}

Publication and peer review are relevant quality signals, not truth labels.
The IOTA and Dirac manuscripts are treated as project evidence; the PSA
semantic-field manuscript is under review. None is counted as independent
replication. Positive, negative, stopped, censored, and withdrawn project
results are retained because a viability audit cannot be built only from
successful gates.

\subsection{Indexed future-law families}

For a candidate retained state or record $m$, intervention family $\Pi$,
observable algebra $\mathcal A$, boundary and resolution declaration $B$,
resource declaration $C$, and horizon $\tau$, define
\begin{equation}
 \Efamily(m;\Pi,\mathcal A,B,C)
 =\left\{P^{\pi}_{\mathcal A}
   \bigl(F_{t:t+\tau}\mid m,B,C\bigr):\pi\in\Pi\right\}.
 \label{eq:future-family}
\end{equation}
The corresponding reachable support for policy $\pi$ is
\begin{equation}
 \Rset^{\pi}_{\tau}(m)=
 \operatorname{supp}P^{\pi}_{\mathcal A}
 \bigl(F_{t:t+\tau}\mid m,B,C\bigr),
 \label{eq:reachable-support}
\end{equation}
with topology and tolerance specified for continuous spaces. A single
continuous trajectory may have probability zero without being impossible.

Equation~\eqref{eq:future-family} is not itself an ICP law. It is shared
comparison language. A \emph{constraint} requires a matched contrast between
families, for example a natural history-formed carrier $m$ and a reference
$m_0$ prepared with the same declared ordinary state variables and resources:
\begin{itemize}
  \item a \textbf{hard constraint} changes reachable support on a common path
  or outcome space;
  \item a \textbf{soft constraint} changes weights while support may remain
  fixed; and
  \item an \textbf{operational constraint} changes a reachability or
  convertibility preorder under the same allowed operations.
\end{itemize}
No universal sign for an entropy difference follows from these definitions.

\subsection{A relative-entropy constraint diagnostic}

For a declared comparison $m$ versus $m_0$ at $x$, one distribution-sensitive member
of the future-family comparison is
\begin{equation}
 \mathfrak C_{x,\tau}^{\pi}
 (m\Vert m_0;\mathcal A,B,C)
 :=D_{\mathrm{KL}}\!\left(
 P^{\pi}_{\mathcal A}(F_{t:t+\tau}\mid m,B,C)
 \middle\Vert
 P^{\pi}_{\mathcal A}(F_{t:t+\tau}\mid m_0,B,C)
 \right),
 \label{eq:icp-constraint-divergence}
\end{equation}
provided both laws use the same outcome space, intervention, horizon, measure,
and coarse-graining, with the first law absolutely continuous with respect to
the second \citep{KullbackLeibler1951Information}. This is an indexed diagnostic,
not the definition of ICP: the full family in Eq.~\eqref{eq:future-family}
retains outcome identity, topology, interventions, and operational order that a
single scalar discards.

The IOTA support quantity has an exact specialization. If the post-intervention
law is uniform on
$\Rset_x^{\mathrm{post}}\subseteq\Rset_x^{\mathrm{prior}}$ and the reference
law is uniform on $\Rset_x^{\mathrm{prior}}$, then
\begin{align}
 D_{\mathrm{KL}}(U_x^{\mathrm{post}}\Vert U_x^{\mathrm{prior}})
 &=\sum_{f\in\Rset_x^{\mathrm{post}}}
   \frac{1}{|\Rset_x^{\mathrm{post}}|}
   \log\frac{|\Rset_x^{\mathrm{prior}}|}
                 {|\Rset_x^{\mathrm{post}}|} \notag\\
 &=\log\frac{|\Rset_x^{\mathrm{prior}}|}
                  {|\Rset_x^{\mathrm{post}}|}
 =\Delta I_x .
 \label{eq:iota-kl-identity}
\end{align}
This direction is essential: reversing the divergence is infinite under strict
support pruning. Non-nested or support-creating interventions require an
explicitly chosen alternative treatment rather than silent reuse of
Eq.~\eqref{eq:iota-kl-identity} \citep{Clark2026IotaConstraint}.

\subsection{The adjacent frontier and signed possibility-space flux}

The adjacent possible becomes operational only after its outcome types and
notion of ``one step'' are declared. Let
$r_\ell:\Omega\rightarrow\mathcal K_\ell$ be a frozen scale-$\ell$ map from
micro-outcomes to registered types or roles, and let $\mathcal V_t^\ell$ be the
types actualized by time $t$. Define
\begin{equation}
 \mathcal J_{t,\tau}^{\pi,\ell}(m_t)
 :=r_\ell\!\left(\Rset_\tau^\pi(m_t)\right)
   \setminus\mathcal V_t^\ell .
 \label{eq:adjacent-frontier}
\end{equation}
When $\tau$ is one registered transition, this is the novelty-bearing sector
of Kauffman's adjacent possible: types not yet actual but reachable in one step
from the actual \citep[p.~142]{Kauffman2000Investigations}. For longer horizons
it is an indexed accessible-novelty set, not a claim of literal one-step
adjacency. In either case it depends on $r_\ell,\pi,\tau,B,$ and $C$; changing a
classifier, resource budget, or policy can change the set.

For matched pre/post supports $R_0$ and $R_1$, retain the identity of the
change rather than only its cardinality:
\begin{equation}
 K=R_0\cap R_1,\qquad
 E=R_1\setminus R_0,\qquad
 L=R_0\setminus R_1,
 \label{eq:possibility-flux-ledger}
\end{equation}
where $E,L,$ and $K$ are enabled, lost, and retained possibilities. A signed
log-volume deformation is
\begin{equation}
 \chi(R_1:R_0):=\log\frac{\mu(R_0)}{\mu(R_1)}.
 \label{eq:signed-support-deformation}
\end{equation}
For nonempty finite uniform nested supports this gives the exact dual pair
\begin{equation}
 \chi(R_1:R_0)=
 \begin{cases}
 D_{\mathrm{KL}}(U_{R_1}\Vert U_{R_0}),&R_1\subseteq R_0
 \quad\text{(pruning)},\\[3pt]
 -D_{\mathrm{KL}}(U_{R_0}\Vert U_{R_1}),&R_0\subseteq R_1
 \quad\text{(enabling)}.
 \end{cases}
 \label{eq:signed-kl-duality}
\end{equation}
Thus $\chi>0$ denotes net contraction and $\chi<0$ net expansion under the
declared measure. If supports are non-nested, $\chi$ is only a signed net-volume
contrast: it can vanish under complete identity turnover. Report
Eq.~\eqref{eq:possibility-flux-ledger} and, where a scalar turnover diagnostic
is needed, a bounded common-support quantity such as
$D_{\mathrm{JS}}(U_{R_1},U_{R_0})$. Empty sets require explicit set reporting
or a declared absorbing outcome; they do not justify arbitrary pseudocounts.

\subsection{The persistent-carrier criterion}

A proposed informational carrier $M$ passes the current ICP criterion only if
four gates are met:
\begin{description}
  \item[Formation.] Controlled histories create the candidate distinction.
  \item[Sufficiency.] The distinction is sufficient for the declared future-law
  family to a frozen error tolerance.
  \item[Carrier intervention.] Ablating, replacing, transplanting, or scrambling
  the proposed carrier changes a preregistered component of the future family,
  while matched shams and resource controls do not.
  \item[Robustness.] The result survives adequate state augmentation and
  reasonable representation, partition, and coarse-graining changes, or its
  representation dependence is stated as a boundary.
\end{description}

Persistence across carrier succession adds one further demand. Let
$c_0,c_1,\ldots$ be physical realizations at successive times. They instantiate
the same role only relative to a declared equivalence relation $\sim_{\Pi,
\mathcal A,B,C,\tau}$ when replacement preserves the relevant future-response
family or a registered relational invariant. Material disjointness alone does
not establish sameness of role; sameness of an analyst's label does not
establish physical transport.

\subsection{Novelty and failure tests}

The strongest null is the complete-state ordinary-dynamics account: once every
causally relevant variable, law, input, boundary, and control is included,
there is no additional historical variable. State augmentation, predictive
causal states, process tensors, control/viability models, and
counterfactual-semantic measures are field-specific versions of this challenge.

ICP earns distinct physical content only by delivering at least one of the
following beyond the best native model:
\begin{enumerate}
  \item a minimal sufficient physically embodied equivalence class;
  \item a representation-stable carrier-succession invariant;
  \item a new impossibility or conversion theorem;
  \item a demonstrably superior compression of intervention-response families;
  or
  \item a novel, preregistered physical prediction.
\end{enumerate}
Repeated disappearance under adequate state augmentation, failure of matched
carrier interventions, dependence on one convenient representation, or the
need to import the target quantum structure would lower or defeat the relevant
claim.
